% Source presentations for register row N-c-019.

\subsection{Five source presentations behind Hoare logic and weakest preconditions}
\label{sec:p09-presentation-n-c-019}

\subsubsection{Floyd's flowchart assertions}

Floyd's program is a sequential flowchart.  Assertions are attached to selected
edges, and each assignment, test, branch, or join produces a local verification
condition.  If the entry assertion holds and every traversed edge preserves the
next assertion, then every execution which reaches the exit satisfies its exit
assertion.  The proof is induction along the finite traversed path
\cite[pp.~19--24]{Floyd1967}.  This is partial correctness: reaching the exit is
a hypothesis of the conclusion.

Termination is a separate construction.  Floyd attaches values in a
well-ordered set and requires strict descent around every relevant cycle.  An
infinite execution would produce an infinite descending sequence, which is
impossible.  The flowchart assertions and the descent values have different
roles; the first do not imply the second \cite[pp.~26--31]{Floyd1967}.  The
countdown derivation in Example~\ref{def:p09-hoare-counter} is the structured-
language descendant of this two-part method, not a claim that Floyd printed a
modern Hoare derivation.

\paragraph{A complete local reconstruction.}
Here is the argument in a small editorial flowchart rather than in modern
triple notation.  Fix an input integer \(n\geq 0\), initialize \(i:=0\) and
\(S:=0\), test \(i<n\), and, on the true edge, execute \(i:=i+1\) followed by
\(S:=S+i\) before returning to the test.  Put on the test edge the assertion

\[
 I\;\equiv\;0\leq i\leq n\ \land\ 2S=i(i+1).
\]

This assertion has one standing only: it is the formula assigned to that
control edge.  Initialization establishes it because both program variables
are zero.  On the true edge, write \(i'=i+1\) and \(S'=S+i'\) for the values
after the two assignments.  From \(I\land i<n\), integrality gives
\(i'=i+1\leq n\), and

\[
 2S'=2S+2(i+1)=i(i+1)+2(i+1)=(i+1)(i+2)=i'(i'+1).
\]

Thus the body returns to the test edge with \(I\) again.  On the false edge,
\(I\land\neg(i<n)\) and \(i\leq n\) give \(i=n\), hence
\(2S=n(n+1)\).  These are exactly the entry, cycle, and exit verification
conditions needed for partial correctness.  Termination is an additional
argument: while the true edge is taken, the natural number \(n-i\) is positive
and the body replaces it by \(n-i-1\).  An infinite traversal would therefore
give an infinite strictly descending sequence of natural numbers.  This
reconstruction exposes, on one program, Floyd's separation between edge-
assertion preservation and well-founded descent
\cite[pp.~24--31]{Floyd1967}.

\subsubsection{Hoare's partial-correctness calculus}

Hoare's 1969 typography is \(P\{Q\}R\): \(P,R\) are assertions and the program
fragment \(Q\) is written inside the braces.  The paper gives assignment,
consequence, composition, conditional, and invariant-based while rules
\cite[pp.~576--578]{Hoare1969}.  The handbook uses the now-standard
\(\{P\}\,Q\,\{R\}\), but the semantics remains partial correctness.

Soundness of the modern registered rules is proved in
Annex~\ref{annex:p9.hoare.1}.  Its proof follows the semantic cases: assignment
evaluates the updated state; sequence supplies the intermediate state;
conditionals split on the guard; and the while case inducts on the finite number
of iterations.  This local proof is not back-attributed as a metatheorem
contained in Hoare's short article.  Hoare also warns that overflow, expression
side effects, range limits, and language-specific parameter behavior can
invalidate mathematical assignment axioms; the exact command language in
Definition~\ref{def:p09-hoare-triple} is therefore part of every theorem.

Hoare's Table~III makes the rule interaction explicit on successive
subtraction.  In the source all variables range over nonnegative integers, and
the program is

\[
 r:=x;\quad q:=0;\quad
 \mathbf{while}\ y\leq r\ \mathbf{do}\ 
      (r:=r-y;\ q:=1+q).
\]

The loop assertion is \(x=r+yq\).  Initialization establishes it in two
assignment steps: \(x=x+y0\), then \(r:=x\), then \(q:=0\).  If
\(x=r+yq\land y\leq r\), the arithmetic equality

\[
 r+yq=(r-y)+y(1+q)
\]

is the substituted precondition for the two body assignments, so composition
restores \(x=r+yq\).  Hoare's iteration rule and one final composition step
therefore derive

\[
 \{\top\}\quad r:=x;\quad q:=0;\quad
 \mathbf{while}\ y\leq r\ \mathbf{do}\ 
      (r:=r-y;\ q:=1+q)\quad
 \{r<y\land x=r+yq\}.
\]

This is the source's twelve-line partial-correctness derivation in compressed
modern typography — not a termination proof
\cite[pp.~578--579, Table~III]{Hoare1969}.  Over the handbook's unbounded
integers, the stronger precondition \(x\geq0\land y>0\), invariant
\(y>0\land r\geq0\land x=r+yq\), and variant \(r\) do yield total correctness:
whenever \(y\leq r\), the next value \(r-y\) is nonnegative and is strictly
smaller than the old \(r\).  That final sentence is an explicitly modern
extension using the registered total-while rule; Hoare states that his 1969
rules themselves provide no basis for successful termination
\cite[p.~579]{Hoare1969}.

\subsubsection{Dijkstra's demonic guarded commands}

Dijkstra's alternative has the form
\[
 \mathbf{if}\ B_1\to S_1\ []\ \cdots\ []\ B_n\to S_n\ \mathbf{fi}.
\]
If several guards are true, any corresponding command may be selected; if none
is true, the alternative aborts.  This command is necessary to express
nondeterministic choice under total correctness.  It is a program command,
differs from the deterministic conditional in the registered while language,
and is used by the source's weakest-precondition equation
\[
 \operatorname{wp}(\mathbf{if}\ \cdots\ \mathbf{fi},R)
 =\left(\bigvee_i B_i\right)
   \land
   \bigwedge_i\bigl(B_i\to\operatorname{wp}(S_i,R)\bigr).
\]
The conjunction over enabled alternatives gives the demonic reading: every
permitted choice must terminate properly in an \(R\)-state
\cite[\S\S2--3.1]{Dijkstra1975}.

Dijkstra first records four laws which constrain this reading:

\[
\begin{aligned}
 \operatorname{wp}(S,\bot)&=\bot,\\
 P\Rightarrow Q&\quad\Longrightarrow\quad
   \operatorname{wp}(S,P)\Rightarrow\operatorname{wp}(S,Q),\\
 \operatorname{wp}(S,P\land Q)
   &=\operatorname{wp}(S,P)\land\operatorname{wp}(S,Q),\\
 \operatorname{wp}(S,P\lor Q)
   &=\operatorname{wp}(S,P)\lor\operatorname{wp}(S,Q)
   \qquad\text{when \(S\) is deterministic.}
\end{aligned}
\]

The final equality is not asserted for a nondeterministic command.  For
example, take
\(R\equiv(m=x\lor m=y)\land m\geq x\land m\geq y\).  Assignment substitution
gives

\[
 \operatorname{wp}(m:=x,R)=x\geq y,
 \qquad
 \operatorname{wp}(m:=y,R)=y\geq x.
\]

Those two weakest preconditions become the guards of Dijkstra's maximum
program.  Their disjunction is valid in a total order, so the alternative
cannot abort; if both hold, either assignment establishes \(R\).  The example
therefore derives the guards from the required result rather than guessing a
deterministic branch order \cite[\S4]{Dijkstra1975}.

The analogous guarded repetition stops only when every guard is false.
Let \(DO\) denote a fixed repetition and \(IF\) the same guarded-command set
enclosed as an alternative, and put \(BB\equiv\bigvee_i B_i\).  Here \(DO\)
and \(IF\) name commands, while \(BB\) is merely the formula saying that at
least one branch is enabled.  For a postcondition \(R\), define predicate
abbreviations

\[
 H_0(R)\equiv R\land\neg BB,
 \qquad
 H_k(R)\equiv
   \operatorname{wp}(IF,H_{k-1}(R))\lor H_0(R)\quad(k>0).
\]

\(H_k(R)\) is the weakest precondition guaranteeing termination in an
\(R\)-state after at most \(k\) selections.  It is neither a new command nor an
asserted invariant; the indexed abbreviation is needed to retain the finite
selection bound, and differs from the unbounded repetition transformer used
immediately in

\[
 \operatorname{wp}(DO,R)\equiv\exists k\geq0\,H_k(R).
\]

The exact invariant--variant hypothesis says that every state satisfying
\(P\land BB\) has \(t\geq0\), and that \(IF\) is guaranteed to terminate from
that state with \(P\) true and its final integer value of \(t\) at most the old
value minus one.  Under that hypothesis, Dijkstra's theorem concludes

\[
 P\Rightarrow\operatorname{wp}(DO,P\land\neg BB).
\]

Here is the finite-approximant proof omitted from the article.  First,
\(H_j(R)\Rightarrow H_{j+1}(R)\) for every \(j\): the case \(j=0\) follows
from the \(H_0(R)\) disjunct in \(H_1(R)\), and the induction step follows
from monotonicity of \(\operatorname{wp}(IF,-)\).  Now use strong induction on
the nonnegative old value \(n\) of \(t\).  From a
\(P\land BB\land t=0\)-state, every outcome of \(IF\) has \(P\) and
\(t'\leq-1\).  Such an outcome cannot still satisfy \(BB\), since the
hypothesis at an enabled \(P\)-state would require \(t'\geq0\).  Every outcome
therefore satisfies \(H_0(P\land\neg BB)\), and the initial state satisfies
\(H_1(P\land\neg BB)\).

For the strong-induction step, start with
\(P\land BB\land t=n\), where \(n>0\).  Every outcome has \(P\) and an integer
\(t'\leq n-1\).  A disabled outcome satisfies \(H_0(P\land\neg BB)\).  At an
enabled outcome the hypothesis gives \(t'\geq0\); the induction hypothesis
places it in \(H_{t'+1}(P\land\neg BB)\), hence in
\(H_n(P\land\neg BB)\) by the monotonicity just proved.  Thus every outcome of
\(IF\) satisfies \(H_n(P\land\neg BB)\), and guaranteed termination gives
\(H_{n+1}(P\land\neg BB)\).  Finally, an arbitrary \(P\)-state is either
already disabled and in \(H_0(P\land\neg BB)\), or is enabled with some
nonnegative integer value \(n\) and lies in
\(H_{n+1}(P\land\neg BB)\).  The defining existential over \(k\) proves the
displayed conclusion.  This reconstruction supplies the induction behind the
source's stated theorem; it does not turn \(\operatorname{wp}\) into the later
weakest liberal precondition \cite[\S3.3]{Dijkstra1975}.

The source's Euclidean derivation supplies all three obligations.  For fixed
positive \(X,Y\), initialize \(x:=X;y:=Y\), maintain

\[
 P\equiv \gcd(X,Y)=\gcd(x,y)\land x>0\land y>0,
 \qquad t\equiv x+y,
\]

and repeat either \(x>y\to x:=x-y\) or
\(y>x\to y:=y-x\).  The equalities
\(\gcd(x,y)=\gcd(x-y,y)\) and
\(\gcd(x,y)=\gcd(x,y-x)\) preserve the gcd component; the guards preserve
positivity; and either subtraction strictly decreases \(t\).  At exit neither
guard holds, so \(x=y\), and \(P\) gives
\(x=y=\gcd(X,Y)\).  This is a total-correctness derivation even though the
choice of an enabled branch may be nondeterministic \cite[\S4]{Dijkstra1975}.

The distinction is visible in one line.  Let
\(C=\mathbf{while}\ \top\ \mathbf{do}\ \mathbf{skip}\).  The partial triple
\(\{\top\}\,C\,\{\bot\}\) is valid because \(C\) never terminates, while
\(\operatorname{wp}(C,\bot)=\bot\) because no initial state guarantees proper
termination.  Thus \(\operatorname{wlp}\) and Dijkstra's \(\operatorname{wp}\)
cannot share one theorem statement.

\subsubsection{Cook's relative completeness}

Cook fixes a deterministic ALGOL-like language with assignments,
conditionals, loops, blocks, and restricted nonrecursive procedures.  His
semantics is partial correctness.  The completeness theorem assumes both that
all valid implications of the assertion theory are available and that the
assertion language can express the strongest post-relations required by the
program constructs \cite[\S\S2--6]{Cook1978}.

The expressibility hypothesis is exact.  Given an assertion \(P\), a command
\(A\), and a tuple \(\vec x\) containing every variable free in \(P\) or in
\(A\), the \emph{post relation} of \(P\) through \(A\) contains precisely those
tuples \(\vec d\) for which some initial state satisfies \(P\), execution of
\(A\) terminates, and the final values of \(\vec x\) are \(\vec d\).  This is a
relation on value tuples, not yet an assertion formula.  The name is necessary
because the completeness induction needs the exact set of reachable outputs,
not an arbitrary sufficient postcondition; it differs from such a
postcondition by requiring both directions of the displayed semantic
description.  Cook calls the assertion language expressive when every one of
these relations is definable by a formula in that language.  The relation is
used immediately as the intermediate assertion for composition
\cite[pp.~85--86]{Cook1978}.

The proof of completeness proceeds by induction after repeatedly replacing
procedure calls by parameter-substituted bodies and renaming local variables
to prevent capture.  This preliminary expansion terminates because Cook's
language excludes recursive procedures.  For a compound
\(\mathbf{begin}\ A;A^*\ \mathbf{end}\), a formula defining the post relation
of \(P\) through \(A\) is the intermediate assertion: both component triples
are semantically true and are shorter induction cases.  For \(x:=e\), the
assignment axiom supplies \(Q[e/x]\), and semantic validity of the original
triple supplies the implication from \(P\) needed by consequence.

The loop case shows why mere first-order consequence completeness is not
enough.  Suppose
\(\{P\}\,\mathbf{while}\ R\ \mathbf{do}\ A\,\{Q\}\) is semantically true, and
let \(\vec y=(y_1,\ldots,y_n)\) list the variables involved.  The required
loop-head relation consists of the states reachable from a \(P\)-state after
some finite number of executions of \(A\), with \(R\) true before each one.
Precisely, let \(\mathcal C_0\) be the relation of \(P\)-state value tuples and
let \(\mathcal C_{k+1}\) contain exactly the final tuples of terminating
executions of \(A\) begun in an \(R\)-state whose tuple is in
\(\mathcal C_k\).  The \(\mathcal C_k\) are metalinguistic relations on value
tuples, not formulas or commands.  They are needed to distinguish the exact
number \(k\) of completed body executions; their union
\(\bigcup_{k\geq0}\mathcal C_k\), used below as the invariant relation, forgets
that exact count.  The proof may not simply add this generally infinitary
union as a disjunction to the assertion language.
Choose fresh variables \(\vec z=(z_1,\ldots,z_n)\), which are not changed by
\(A\), and form the auxiliary command

\[
 W_{\vec z}\equiv
 \mathbf{while}\!
 \left(R\land\bigvee_{i=1}^{n}y_i\neq z_i\right)
 \mathbf{do}\ A.
\]

This notation is eliminable: \(W_{\vec z}\) abbreviates only the displayed
while command and is used once in the construction.  If
\(\Phi(\vec y,\vec z)\) is a formula expressing the post relation of \(P\)
through \(W_{\vec z}\), then

\[
 J(\vec y)\equiv\exists z_1\cdots\exists z_n\,
                    \Phi(\vec y,\vec z)
\]

defines exactly the finite-reachability relation above.  Here \(J\) is an
assertion formula, not a new program construct: choosing \(\vec z\) equal to a
desired reachable loop-head tuple makes the auxiliary loop stop there, while
every terminating run of the auxiliary loop consists only of genuine body
iterations.  The zero-iteration case proves \(P\Rightarrow J\); one more body
step proves \(\{J\land R\}\,A\,\{J\}\); and semantic truth of the original
triple proves \(J\land\neg R\Rightarrow Q\).  The while and consequence rules
then finish this induction case \cite[pp.~86--87]{Cook1978}.

Procedure calls require the remaining source rules: fresh renamings remove
formal/actual clashes, parameter substitution instantiates the proved body
property, and additional fresh variables retain values which must survive the
substitution.  Those restrictions are part of the theorem's language, not
administrative decoration \cite[pp.~87--89]{Cook1978}.

The result is relative because neither the truth of the assertion theory nor
the required expressibility is supplied by the Hoare rules.  Cook makes the
first obstruction precise: if the assertion proof system is axiomatic, its
derivable triples are recursively enumerable, whereas
\(\{\top\}\,A\,\{\bot\}\) is true exactly when \(A\) fails to halt for every
initial assignment to its global variables.  In a language with an
undecidable halting problem, all true triples therefore cannot be recursively
enumerable.  Theorem~\ref{thm:p09-hoare-relative-completeness} does not claim a
decision procedure, total correctness, or completeness for a different
language \cite[pp.~85--86]{Cook1978}.

\subsubsection{Apt's language-fragment calibration}

Apt's 1981 survey distinguishes partial from total correctness and tracks how
soundness and completeness change when the command language gains procedures,
local variables, and different parameter mechanisms \cite{Apt1981}.  Its role
is comparative: it does not make every earlier rule or proof an Apt result.

For total correctness of a while loop, Apt prints a rule which the survey
explicitly says was formulated by Harel.  Let \(p(n)\) be an assertion whose
free variable \(n\) ranges over the natural numbers and does not occur in the
body \(S\).  Every triple in this paragraph has total-correctness standing.
The rule is

\par\noindent
\begin{minipage}{\linewidth}
\showcalculus{n-c-019-harel-indexed-while}
\end{minipage}
\par

The indexed assertion family has an exact standing here: \(p(n)\) is an
assertion formula parameterized by a rigid natural-number value, not a family
of commands or a new semantic relation.  The index is necessary to count the
remaining body executions; it differs from an ordinary invariant because each
body execution must lower the index from \(n+1\) to \(n\), and the rule uses
that distinction in its middle premise.  This is not the
decreasing-integer-variant rule on the handbook sheet.

Apt applies the indexed rule to the same successive-subtraction division
program.  Write

\[
 r\equiv a\,y+b=x\land b\geq0,
 \qquad
 p(n)\equiv r\land ny\leq b<(n+1)y .
\]

Under the standard integer interpretation and \(y>0\), \(r\) implies
\(\exists n\,p(n)\).  From \(p(n+1)\) one has \(b\geq y\); after
\(b:=b-y;a:=a+1\), the equality \(a\,y+b=x\) is preserved and
\(ny\leq b<(n+1)y\), so \(p(n)\) holds.  Finally \(p(0)\) gives
\(0\leq b<y\), hence the guard \(b\geq y\) is false.  These are exactly the
three premises of the printed rule
\cite[pp.~441--442]{Apt1981}.

The interpretation restriction matters.  Apt observes that
\(\mathbf{while}\ x>0\ \mathbf{do}\ x:=x-1\) terminates over the standard
natural numbers but need not terminate in a nonstandard model of arithmetic.
Consequently, no proof system can both be totally sound over all
interpretations and prove even this standard countdown from the full set of
standard arithmetic truths.  The subsequent soundness and completeness result
therefore restricts attention to arithmetical interpretations containing a
standard natural-number part and, for completeness, a definable coding of
finite sequences \cite[pp.~442--443]{Apt1981}.  Elsewhere in the survey,
auxiliary variables have the separate purpose of retaining program-variable
values across procedure calls; that substitution discipline must not be
confused with the natural index in the total-while rule
\cite[p.~448]{Apt1981}.

\begin{table}[htbp]
\centering\scriptsize
\begin{tabularx}{\textwidth}{@{}l X X X X@{}}
\toprule
Source & Command language & Termination standing & Assertion basis &
Source status and exact role \\
\midrule
Floyd 1967 & sequential flowcharts; formulas label selected edges &
partial correctness; separate well-ordered descent & assertions over machine
states at control points & local preservation and descent method with worked
flowcharts \\
Hoare 1969 & simple structured sequential commands and triples &
partial while rule only & mathematical assertions, valid implications, and
chosen arithmetic axioms & rule schemas and a twelve-line worked derivation;
no general metatheoretic soundness proof \\
Dijkstra 1975 & demonic guarded alternatives and repetitions &
proper termination is part of \(\operatorname{wp}\) & predicates on the full
state space satisfying the transformer laws & exact equations and stated
repetition theorems, followed by maximum and gcd derivations \\
Cook 1978 & deterministic ALGOL-like fragment with nonrecursive procedures &
partial correctness & semantically complete assertion theory plus definability
of every required post relation & formal rule soundness and relative-
completeness proofs \\
Apt 1981 & separately surveyed while, procedure, local-variable, and parameter
fragments & partial and total kept separate & fragment-specific interpretation
hypotheses; arithmetical interpretations for the displayed total system &
survey proofs and proof sketches, with Harel attribution for while Rule~II \\
\bottomrule
\end{tabularx}
\caption{Noninterchangeable source roles for Floyd--Hoare reasoning.}
\label{tab:p09-hoare-source-roles}
\end{table}
